\section{Motivation}\label{sec:motivation}

How much useful mathematical work can an AI research agent produce when
given an existing problem collection, mathematical software and two days?
The task includes recognizing prior solutions, choosing approachable
questions, finding arguments and checking their scope. It is closer to
research than to a fixed set of exercises, and its output requires a
different assessment: a plausible proof and a successful computation are
evidence to examine, not automatic credit for a new theorem.

Following our Kourovka experiment~\cite{kourovka-paper,kourovka-repo},
Vladimir Shpilrain suggested using the GroupWorld collection that he
hosts~\cite{groupworld}. The collection was originally selected in
1997--99 by Gilbert Baumslag, Alexei Myasnikov and Shpilrain and includes
questions about free, nilpotent, metabelian and hyperbolic groups, braids,
equations and algorithms. Its background notes and Hall of Fame are part
of the input, rather than material to ignore when counting successes.
The preparation archive has 195 numbered entries in 16 categories;
several have multiple parts, and many already have known answers.

The experiment ran from 28 September 2026, 10:04:49 UTC, to
30 September 2026, 10:04:49 UTC, with requested local limits of generally
20 CPU cores and 100 GB RAM. These explicit limits give the exercise a
competition-like format, but the questions were drawn from a historical
research collection rather than a freshly curated test set. The agent
could use the literature, GAP, exact programs and formal proof tools.
There was no restriction to producing a particular proof style or to
using a theorem prover for the final argument.

The mathematical output is varied. An explicit virtual isomorphism of
an index-two subgroup of a free group is proposed to have no nontrivial
forward-invariant subgroup. A short family of virtually free presentations
forces superpolynomial explicit Dehn output. Other arguments propose
decisions for nilpotent direct decomposition and single commutators,
and use finite-field Fox calculus, formal languages, model theory and
tree actions. Two partial results give a computable growth constant with
a certified interval and an infinite special-braid family. Section~\ref{sec:inventory}
lists every counted entry before the proofs.

This version is an account of the deadline candidates, prepared after
the run. It does not report an independent mathematical referee's verdict.
Correctness and novelty remain separate questions for each argument.
The earlier Kourovka paper includes a second model's review and later
human corrections, including the withdrawal of an argument applied to
the wrong field. That experience motivated the separate statement,
proof and attribution checks here; it does not validate this portfolio.
The two runs have different inputs and review histories, so their raw
candidate counts are not a controlled comparison of performance.

The rapidly growing work on language models, learning-guided proving
and AI for mathematics precludes a proper survey here. We cite the
mathematical ingredients actually used and the directly preceding
experiment. Computational details and the larger proof interfaces are
placed in the appendices; readers may start with the explicit arguments
in Section~\ref{sec:explicit} or the subjects that interest them.
